\section{Threats to validity}
\label{sec:threats}

\paragraph{We inspected the implementation, and then executed it} Two
limitations are easily conflated, so we separate them.

\emph{Source.} The realisation of \mbox{MST-wi} is
public~\cite{Chaleshtari2022Metamorphic}. We obtained it, read the predicates at
source (Section~\ref{sec:impl}), and built it under JDK~8 and Maven together with
its \ResArtAuthzMrTotal{} authorization relations, which load. Where we describe a
predicate's semantics we cite the code rather than the prose.

\emph{Execution.} We then ran those relations---their source unmodified---against
a real system; the only additions were the build configuration
(\texttt{pom.xml}) and an external harness/driver that the upstream artifact
does not ship
(Section~\ref{sec:native-run}). This was not free: the library does not ship the
per-system configuration and crawl data that its own demonstrations reference;
\ifnum\ResArtFactFSixupstreampomdeclaresmrsourceroot=0
nor does its build configuration compile the relations themselves, which live in a
source directory the build file does not declare---a declaration we had to add,
and Section~\ref{sec:native-run} reports that one-line diff along with the rest
of our footprint;\fi{}
executing a relation therefore required assembling a subject system, a system
configuration and a crawl, which is what we did. The harness is ours, and its
cost must be stated rather than assumed away: we execute without a browser,
substituting the HTTP observation channel while leaving the relations, the
predicates and the loaded crawl records untouched. That separation is exact for
the reachability predicate, which is set-membership over loaded records and issues
no request; it is \emph{not} exact for the content predicate, which consumes the
observed output, and that predicate happens not to be exercised by our data. We
report its invocation count rather than claim a cleanliness we did not
demonstrate.

\emph{What the run buys, and what it does not.} It converts one class of claim:
we can now say that a false proof is produced by the \emph{published} relations
under the \emph{published} predicates on a real system, and not only by our
reading of them. It does not convert the rest. The \ResSynRecords{}-record
factorial and the cell study remain statements about our instantiation, because
the harness is ours and because our exception vocabulary is drawn from the same
predicate family: where we approximate an exception semantics we say so, and our
result there is a statement about that family's expressive limit, not a
measurement of the engine's false-positive rate. We still do not report a
sensitivity or a specificity attributable to \mbox{MST-wi}, and we still do not
assert that it fails.

\paragraph{Metric convention} The specificity of
Section~\ref{sec:synthetic-ablation} uses two follow-up inputs per case as a
stand-in for the engine's input count. That approximation is not the engine's
number, and we use it only to exhibit the denominator's resolution. The
conclusion we draw---that an input-unit denominator cannot see a
scenario-uniform failure---does not depend on the exact count.

\paragraph{The independent unit is the design cell, not the record} As set out
in Section~\ref{sec:stats}, the \ResSynRecords{} records expand
\ResDesignCells{} cells by \ResDesignMrs{} MRs and by two switches that we throw
deliberately. Records inside a cell are not independent. All intervals in this
paper are consequently conventional and optimistic, which makes our
\emph{negative} finding---non-significance of the recall difference in
Boundary~2---conservative rather than fragile: a narrower interval only makes a
difference easier to call significant, and we still cannot call ours.

\paragraph{Two MRs, and a catalog of one family} We instantiated
\ResDesignMrs{} MRs, both from the same family---authorization relations whose
gate is the absence of evidence. The native run of Section~\ref{sec:native-run}
executes \NtMrExecuted{} of the catalog's \ResArtAuthzMrTotal{} authorization
relations, which lifts the count but not the family: \ResArtAuthzGatedCrtg{} of
them gate on the reachability predicate and \ResArtAuthzGatedUcrc{} on the
content predicate, and only the former produce alarms in our run, so the boundary
we report remains a property of that gate. What the run
broadens is the illustration of the boundary, not the evidence that the boundary
is where we say. The ceiling on any generalization to ``\mbox{MST-wi}'s MRs''
therefore stands, and we do not generalize. This remains the single most
important limitation of the paper, and it is the one we would address first with
more work: raising it requires relations whose gate is \emph{not} of this family,
which is a different experiment from the one reported here.

\paragraph{Second implementation, shared blind spot} The second implementation
agreed with the first on \ResAltAgree{} of \ResAltCompared{} records
(\ResAltAgreeRate{}; Wilson 95\% CI
$[\ResAltAgreeRateLo{}, \ResAltAgreeRateHi{}]$). Both were written from the same
reading of the predicate semantics, so this bounds implementation accident only.
It does not bound a shared misreading, and we do not present it as
corroboration.

\paragraph{Real systems, all correct, none vulnerable} The real-system study
spans three systems---Gitea \ResRsVersion{}, Gogs \GgVersion{} (lineage-related:
Gitea forked from Gogs in 2016), and GitLab \GlVersion{}
(no lineage)---chosen for accessibility rather than for representativeness, and
none contains an object-level authorization vulnerability. The experiment
therefore measures only the false-proof side and cannot speak to detection on a
vulnerable real system. The replication is across git forges, one domain; it is
not a sample of web applications at large.

\paragraph{Channel proportions are a design parameter} The breakdown of
Table~\ref{tab:rs-m6}---including the \ResRsChannelNotObservableRate{} share of
authorized access that the framework's log cannot see---reflects the topology we
constructed to exercise the mechanism. It is not a field prevalence, and it must
not be quoted as one. The instance-independent measurements in this paper are the
grant-endpoint census (Table~\ref{tab:rs-m1}) and the predicate counts over the
public artifact (Section~\ref{sec:background}).

\paragraph{Our own earlier claims are among the casualties} An earlier iteration
of this work proposed the state-calibrated adjudicator v3 as a mechanism and
reported its precision of \ResAbVThreePrec{}. Two things were wrong with that.
The precision was bought with \ResAbVThreeFN{} false negatives, so its F1 is
\ResAbVThreeFone{} and quoting the precision alone was a selective presentation;
and the control that we had used to argue for the mechanism's necessity
(\texttt{direct\_only}, \ResAbDirectOnlyPrec{} precision) was a straw man because
it discarded the access result. Replacing it with a fair control
(\texttt{direct\_observe}) is what produced the negative result of
Section~\ref{sec:synthetic-ablation}. We record this because it defines the
condition under which our own claim would have been wrong, and a reader should
demand that condition of any similar claim: a control that consumes the same
observable quantity as the true label.

\paragraph{Constructed scenarios} The synthetic system is ours, its scenarios
are designed to separate causes, and its ground truth is defined by
construction. Section~\ref{sec:realsystem} exists precisely to test whether the
result survives outside that construction, and it does, in the direction and
with the endogenous control we report; but three systems, two of them related by
lineage, tested once each, are still not a population.
